% Part: sets-functions-relations
% Chapter: sets
% Section: pairs-and-products
% Hindi translation (hi-Deva-IN) of Open Logic commit 9620cc73f9c8e0ad003c514a5d3748f29611c4c0.

\documentclass[../../../include/open-logic-section]{subfiles}

\begin{document}

\olfileid[hi]{sfr}{set}{pai}
\olsection{युग्म, क्रमित अवयव-समूह और कार्तीय गुणनफल}

\begin{explain}
विस्तारिता से यह निष्कर्ष निकलता है कि किसी समुच्चय के अवयवों का कोई
क्रम नहीं होता। इसलिए यदि हमें क्रम निरूपित करना हो, तो हम
\emph{क्रमित युग्म} $\tuple{x, y}$ का उपयोग करते हैं। अक्रमित युग्म
$\{x, y\}$ में क्रम का कोई महत्त्व नहीं होता:
$\{x, y\} = \{y, x\}$। क्रमित युग्म में क्रम महत्त्वपूर्ण होता है: यदि $x
\neq y$, तो $\tuple{x, y} \neq \tuple{y, x}$।

समुच्चय सिद्धान्त में क्रमित युग्मों को हमें किस प्रकार समझना चाहिए?
मुख्य बात यह है कि हम इस विचार को बनाए रखना चाहते हैं कि दो क्रमित युग्म
तभी और केवल तभी समान हों, जब उनके प्रथम अवयव समान हों और उनके द्वितीय
अवयव भी समान हों; अर्थात्:
\[
  \tuple{a, b}= \tuple{c, d}\text{ तभी और केवल तभी जब }a = c \text{ और }b=d.
\]
समुच्चय सिद्धान्त में हम वीनर--कुरातोव्स्की परिभाषा की सहायता से क्रमित
युग्मों को परिभाषित कर सकते हैं।
\end{explain}

\begin{defn}[क्रमित युग्म]\ollabel{wienerkuratowski}
	$\tuple{a, b} = \{\{a\}, \{a, b\}\}$.
\end{defn}

\begin{prob}
	\olref[sfr][set][pai]{wienerkuratowski} का उपयोग करके सिद्ध कीजिए कि
	$\tuple{a,
	b}= \tuple{c, d}$ तभी और केवल तभी, जब $a = c$ और $b=d$
	दोनों हों।
\end{prob}

\begin{explain}
क्रमित युग्म की परिभाषा निर्धारित कर लेने के बाद हम उससे और समुच्चय
परिभाषित कर सकते हैं। उदाहरण के लिए, कभी-कभी हमें दो से अधिक वस्तुओं के
क्रमित अनुक्रम भी चाहिए होते हैं, जैसे \emph{त्रिक} $\tuple{x, y, z}$,
\emph{चतुष्टय} $\tuple{x, y, z, u}$, इत्यादि। हम त्रिक को ऐसा विशेष
क्रमित युग्म मान सकते हैं जिसका प्रथम अवयव स्वयं एक क्रमित युग्म हो:
$\tuple{x, y, z}$ का अर्थ $\tuple{\tuple{x, y},z}$ है। चतुष्टय के लिए
भी यही बात लागू होती है: $\tuple{x,y,z,u}$ का अर्थ
$\tuple{\tuple{\tuple{x,y},z},u}$ है, और इसी प्रकार आगे भी। सामान्यतः
हम \emph{क्रमित $n$-अवयवी समूह} $\tuple{x_1, \dots, x_n}$ की बात करते
हैं।

क्रमित युग्मों अथवा अन्य क्रमित $n$-अवयवी समूहों के कुछ समुच्चय हमारे
लिए उपयोगी होंगे।
\end{explain}

\begin{defn}[कार्तीय गुणनफल]
दिए गए समुच्चयों $A$ और $B$ का \emph{कार्तीय गुणनफल} $A \times B$ इस
प्रकार परिभाषित होता है:
\[
  A \times B = \Setabs{\tuple{x, y}}{x \in A \text{ और } y \in B}.
\]
\end{defn}

\begin{ex}
यदि $A = \{0, 1\}$ और $B = \{1, a, b\}$, तो उनका गुणनफल है
\[
A \times B = \{ \tuple{0, 1}, \tuple{0, a}, \tuple{0, b},
    \tuple{1, 1}, \tuple{1, a}, \tuple{1, b} \}.
\]
\end{ex}

\begin{ex}
यदि $A$ एक समुच्चय है, तो स्वयं $A$ के साथ उसका गुणनफल $A \times A$,
$A^2$ भी लिखा जाता है। यह ऐसे \emph{सभी} युग्मों $\tuple{x, y}$ का
समुच्चय है जिनके लिए $x, y \in A$। सभी त्रिकों $\tuple{x, y, z}$ का
समुच्चय $A^3$ है, और इसी प्रकार आगे भी। हम इसकी पुनरावर्ती परिभाषा दे
सकते हैं:
\begin{align*}
  A^1 & = A\\
  A^{k+1} & = A^k \times A
\end{align*}
\end{ex}

\begin{prob}
$\{1, 2, 3\}^3$ के सभी !!{element}ों की सूची दीजिए।
\end{prob}

\begin{prop}\ollabel{cardnmprod}
यदि $A$ में $n$ !!{element} और $B$ में $m$ !!{element} हैं, तो $A
\times B$ में $n\cdot m$ अवयव हैं।
\end{prop}

\begin{proof}
प्रत्येक ऐसे !!{element}~$x$ के लिए जो~$A$ में है, $m$ !!{element}
$\tuple{x, y} \in A \times B$ रूप के होते हैं। मान लीजिए
$B_x = \Setabs{\tuple{x, y}}{y
  \in B}$। जब भी $x_1 \neq x_2$ होता है, तब
$\tuple{x_1, y} \neq
\tuple{x_2, y}$ होता है; इसलिए $B_{x_1} \cap B_{x_2} = \emptyset$। किंतु
यदि $A = \{x_1,
\dots, x_n\}$, तो $A \times B = B_{x_1} \cup \dots \cup B_{x_n}$;
अतः इसमें $n\cdot m$ !!{element} होते हैं।

इसे दृष्टिगोचर बनाने के लिए~$A \times B$ के !!{element}ों को एक जालिका
में व्यवस्थित कीजिए:
\[
\begin{array}{rcccc}
  B_{x_1} = & \{\tuple{x_1, y_1} & \tuple{x_1, y_2} & \dots & \tuple{x_1, y_m}\}\\
  B_{x_2} = & \{\tuple{x_2, y_1} & \tuple{x_2, y_2} & \dots & \tuple{x_2, y_m}\}\\
  \vdots & & \vdots\\
  B_{x_n} = & \{\tuple{x_n, y_1} & \tuple{x_n, y_2} & \dots & \tuple{x_n, y_m}\}
\end{array}
\]
चूँकि सभी $x_i$ परस्पर भिन्न हैं और सभी $y_j$ भी परस्पर भिन्न हैं, इसलिए
इस जालिका के कोई भी दो युग्म समान नहीं हैं और इसमें कुल $n\cdot m$ युग्म
हैं।
\end{proof}

\begin{prob}
$k$ पर आगमन द्वारा सिद्ध कीजिए कि प्रत्येक $k \ge 1$ के लिए, यदि $A$
में $n$ !!{element} हैं, तो $A^k$ में $n^k$ !!{element} हैं।
\end{prob}

\begin{ex}
यदि $A$ एक समुच्चय है, तो $A$ के !!{element}ों के किसी भी अनुक्रम को
$A$ पर बना एक \emph{शब्द} कहते हैं। किसी अनुक्रम को ऐसे $n$-अवयवी समूह
के रूप में समझा जा सकता है जिसके !!{element},~$A$ से लिए गए हों। उदाहरण के लिए, यदि
$A = \{a, b, c\}$, तो अनुक्रम ``$bac$'' को त्रिक
$\tuple{b, a, c}$ माना जा सकता है। शब्द, अर्थात् प्रतीकों के अनुक्रम,
कंप्यूटर विज्ञान में अत्यंत महत्त्वपूर्ण हैं। परिपाटी के अनुसार हम $A$
के !!{element}ों को लंबाई~$1$ के अनुक्रम और $\emptyset$ को लंबाई~$0$
का अनुक्रम मानते हैं। तब $A$ पर बने \emph{सभी} शब्दों का समुच्चय है
\[
A^* = \{\emptyset\} \cup A \cup A^2 \cup A^3 \cup \dots
\]
\end{ex}

\end{document}
